\subsection{Zuordnung eines neuen Personenprofils}
\label{sec:new-profile}

Solange eine Spur noch keine Personenkennung besitzt, wird sie als Kandidat für ein neues Profil betrachtet. Die Personenausschnitte werden dabei aus dem ursprünglichen Frame entnommen, \textbf{bevor} Boxen oder Trackkennungen eingezeichnet werden. Das System prüft zunächst für jeden Ausschnitt der Spur, ob er geometrisch gültig ist und die Mindestgröße erfüllt.

Anschließend ermittelt das System Schärfe, Helligkeit, Grö{\ss}e, Seitenverhältnis, Randkontakt und Detektionskonfidenz. Daraus wird ein gewichteter Qualitätswert $Q$ zwischen null und eins berechnet. Schärfe und Grö{\ss}e erhalten dabei das höchste Gewicht.

Für Initialkandidaten gelten zusätzlich eigene Mindestgrenzen, die unabhängig vom Gesamtwert $Q$ geprüft werden: Die Schärfe muss mindestens $s_{\mathrm{init}}=0{,}40$ erreichen, das Seitenverhältnis mindestens $a_{\mathrm{init}}$. Berührt die ursprüngliche Bounding Box einen Bildrand, gilt zusätzlich die strengere Schärfegrenze $s_{\mathrm{Rand}}=0{,}45$; gleichzeitig sinkt der Randfaktor $R$ in der Berechnung von $Q$ auf $0{,}65$. Ein randberührender, unscharfer Ausschnitt wird damit doppelt erfasst: einmal über den niedrigeren Gesamtwert $Q$, einmal über die eigene Schärfegrenze. Eine gute Grö{\ss}e oder Helligkeit kann einen offensichtlich unscharfen oder untypisch breiten Ausschnitt folglich nicht allein über $Q$ ausgleichen. Details zu diesen Grenzen stehen in Anhang~\ref{app:profile-details}.

Erst wenn alle Grenzen erfüllt sind, berechnet OSNet-x1.0 ein 512-dimensionales Embedding~\cite{zhou2019osnet}, dessen Gewichte aus dem Training auf MSMT17 stammen~\cite{torchreid-modelzoo}. Eigenes Training oder Fine-Tuning findet nicht statt.

Vor der ersten Zuordnung sammelt das System fünf akzeptierte Ausschnitte derselben Spur. Zwischen zwei angenommenen Ausschnitten müssen dabei mindestens drei Frames liegen, damit nicht nur nahezu identische Nachbarbilder in das Initialprofil eingehen; dieser Abstand wird ausschließlich zwischen angenommenen Ausschnitten gezählt, ein abgelehnter Versuch verändert ihn nicht. Ein abgelehnter Ausschnitt füllt den Initialpuffer entsprechend nicht, und die Spur bleibt so lange unbekannt, bis eine spätere, besser geeignete Beobachtung vorliegt.

Sobald fünf Ausschnitte angenommen wurden, werden ihre normalisierten Embeddings $z_i$ qualitätsgewichtet zusammengeführt:
\begin{equation}
\bar z=\operatorname{norm}\left(\frac{\sum_i w_i z_i}{\sum_i w_i}\right),
\qquad w_i=\max(Q_i,0{,}05).
\label{eq:fusion}
\end{equation}
Das feste Mindestgewicht $0{,}05$ sorgt dafür, dass auch bei deaktivierten Qualitätsschwellen jeder Beitrag zählt.

Gegen dieses zusammengeführte Embedding vergleicht das System alle zulässigen gespeicherten Profile mittels Cosine Similarity. Für zwei normalisierte Vektoren $z_q$ und $z_p$ entspricht dies dem Skalarprodukt:
\begin{equation}
s(z_q,z_p)=z_q^\top z_p.
\label{eq:cosine}
\end{equation}
Bereits sichtbare Personenkennungen werden dabei von der Auswahl ausgeschlossen, um Doppelzuweisungen zu vermeiden. Erreicht das ähnlichste Profil den konfigurierten Ähnlichkeitsschwellwert $\tau$, wird seine Personenkennung der Spur zugeordnet. Andernfalls wird ein neues Profil mit einer neuen Personenkennung angelegt und der Spur zugeordnet. Neben dem Ergebnis werden auch das ähnlichste geprüfte Profil, sein Score und der Ablehnungsgrund protokolliert. Dadurch bleibt ein knapp verfehlter Profilabgleich nachträglich nachvollziehbar.

Im ersteren Fall wird das zusammengeführte Embedding zusätzlich gegen die Update-Ähnlichkeitsschwelle aus Abschnitt~\ref{sec:profile-update} geprüft, bevor es in das bestehende Profil einfließt. Eine wiedergefundene Person durchläuft somit denselben Update-Schutz wie ein reguläres Update. Wird diese Schwelle nicht erreicht, bleibt die Personenkennung zugeordnet, das Profil selbst aber unverändert.